\documentclass[12pt]{article}

% La gestion des fontes, sachant que ce document est sensé
% être compilé au moyen de XeLaTeX
%\usepackage{fontspec,xltxtra}

% Pour permettre la redéfinition des dimensions 
\usepackage[a4paper]{geometry}

% Pour le support du français
\usepackage[francais]{babel}

% Les polices à utiliser
%\newcommand{\fontsphil}{
 % \defaultfontfeatures{Mapping=tex-text}
  %\defaultfontfeatures{Ligatures=TeX}
  %\setmainfont{Linux Libertine O}
  %\setsansfont{Liberation Sans}
  %\setmonofont[Scale=.9]{Ubuntu Mono}
  %\newfontfamily{\hwfont}{Comic Sans MS}
  %\newfontfamily{\fatfont}{Arial Black}
  %\newfontfamily{\authorfont}{Liberation Serif}
  %\newfontfamily{\spacefont}{Liberation Mono}
  %\fontspec[SmallCapsFont={Arial Black}]{Arial Black}
%}
%\fontsphil

% Des polices pour quelques caractères spéciaux
\usepackage{wasysym}
\usepackage{pifont}

% Gestion des images
\usepackage{graphicx}

% Gestion des couleurs
\usepackage{color}
\definecolor{fondsection}{rgb}{.9,.9,.9}

% Gestion des espaces
\usepackage{xspace}

% Gestion des URL et liens hypertexte
\usepackage{hyperref}

% Configuration des dimensions
\geometry{%
  left=20mm,width=165mm,
  top=15mm,height=267mm,
  footskip=15mm}

% Le titre
\title{Mon super rapport de stage}
\author{Alain Térieur}
\date{Année universitaire 2022--2023}

% Une macro
\newcommand{\oula}{(\textbf{???})}

% Macro pour page vide sans numérotation
\usepackage{afterpage}
\newcommand\myemptypage{
    \null
    \thispagestyle{empty}
    \addtocounter{page}{-1}
    \newpage
    }

% Packet pour les tableaux
\usepackage{array}
\usepackage{tabularx}

%======================================================================
% D É B U T   D U   D O C U M E N T
%======================================================================

\begin{document}

\thispagestyle{empty}
\begin{center}
  \includegraphics[width=12cm]{Logo_IUT-UL_dept_info.png}
\end{center}

\vspace{1cm}

\noindent
{\large
  IUT Nancy Charlemagne\\
  Université de Lorraine\\
  2 ter boulevard Charlemagne\\
  BP 55227\\
  54052 Nancy Cedex\\[5mm]
  Département informatique
}

\vspace{5cm}

\begin{center}
  {\huge
    \textbf{Synthèse des TD}
  }
\end{center}

\vspace{5cm}

% \vspace{2cm}
\vfill

{\Large
  \noindent
  Etudiant : Théo Eisenbart\\
  Année universitaire 2022--2023
}
% Ajout d'une page vide
\newpage
\thispagestyle{empty}
\mbox{}
\newpage

\newpage
% Table des matières
\tableofcontents

\newpage

\section{Efficacité de l'environnement de travail}}

	\subsection{TD-1 : La souris}

\begin{itemize}
	\item Désactiver votre souris au noveau système avec la commande xinput
\end{itemize}

\vspace{0.3cm}
\$ sudo xinput disable\{2,4,6\}
\vspace{0.3cm}

\begin{itemize}
	\item Initialiser un fichier dans lequel nous allons lister tous les problèmes d'efficacité rencontrés pendant cette séance
\end{itemize}
\vspace{0.3cm}

\begin{tabular}{|c|p{5cm}|p{10cm}|}
	\hline
	\textbf{Priorité} & \textbf{Problème} & \textbf{Correctif}\\ 
	\hline
	1 & Logout difficile au clavier & Raccourci clavier \textbf{Ctrl+Alt+Suppr/Delete}\\
	\hline
	2 & Impossible d'éditer des documents PDF avec Google Drive & Utilisation de LaTeX\\
	\hline
	3 & Impossible d'avoir plusieurs terminaux en parallèle & Sous Terminator, \textbf{Ctrl+Shift+O} pour split le terminal verticalement,\newline \textbf{Ctrl+Shift+L} pour split le terminal horizontalement\\ 
	\hline
	4 & Naviguer dans discord sans la souris & \textbf{Tab} pour se déplacer de haut en bas, \textbf{Shift+Tab} pour aller de bas en haut \newline et \textbf{Ctrl+Tab} pour naviguer de gauche à droite\\
	\hline 
	5 & Se déplacer dans un éditeur de text sans la souris & Utilisation de Vim\\
	\hline
	6 & Ouvrir les settings de notre environnement graphique & shortcut configurable dans les settings ex : \textbf{Alt+S}\\
	\hline
	7 & Accèder au navigateur de fichier & shotcut configurable dans les settings ex : \textbf{Alt+F}\\
	\hline
\end{tabular}

\vspace{0.3cm}
	\subsection{TD-2 : Le clavier}
\begin{itemize}
	\item Identifier un site permettant de s'améliorer à taper au clavier
\end{itemize}

\end{document}
